% !TEX root = ../../main.tex
%% Corrigé du défi du chapitre 4 — vérifié par verifications/defis_01_05.py.

\begin{blocLivre}{cor-4-25a}
\begin{corrige}{25}{Oscillateur amorti, régime forcé et résonance}
\demarche{Équation caractéristique pour le régime libre, exponentielle
complexe pour le régime forcé, étude d'un trinôme pour la résonance.}
\partieExo{A}{Le régime libre ($f=0$)}
\begin{questions}
  \item Équation caractéristique $r^{2}+2\lambda r+\omega_0^{2}=0$, de
    discriminant réduit $\lambda^{2}-\omega_0^{2}$.
    \begin{itemize}[label={\color{otNavy}\textbullet},leftmargin=1.2em,itemsep=0.4mm,topsep=0.6mm]
      \item $\lambda>\omega_0$ (apériodique) : deux racines réelles
        $r_\pm=-\lambda\pm\sqrt{\lambda^{2}-\omega_0^{2}}$,
        $y=C_1\mathrm{e}^{r_+t}+C_2\mathrm{e}^{r_-t}$ ;
      \item $\lambda=\omega_0$ (critique) : racine double $-\lambda$,
        $y=(C_1+C_2t)\,\mathrm{e}^{-\lambda t}$ ;
      \item $\lambda<\omega_0$ (pseudo-périodique) : racines
        $-\lambda\pm\mathrm{i}\Omega$,
        $y=\mathrm{e}^{-\lambda t}\bigl(C_1\cos\Omega t+C_2\sin\Omega t\bigr)$.
    \end{itemize}
  \item $(2,\sqrt3)$ : racines $-1$, $-3$ ; $C_1+C_2=1$ et
    $-C_1-3C_2=0$ donnent
    $\res{y=\tfrac32\mathrm{e}^{-t}-\tfrac12\mathrm{e}^{-3t}}$.
    $(1,1)$ : $\res{y=(1+t)\,\mathrm{e}^{-t}}$.
    $(1,2)$ : $\Omega=\sqrt3$, $C_1=1$, $-C_1+\sqrt3\,C_2=0$, d'où
    $\res{y=\mathrm{e}^{-t}\bigl(\cos\sqrt3\,t+\tfrac{1}{\sqrt3}\sin\sqrt3\,t\bigr)}$.
  \item Dans le cas apériodique, $r_-<r_+<0$ car
    $\sqrt{\lambda^{2}-\omega_0^{2}}<\lambda$ ; dans les deux autres,
    $t^{k}\mathrm{e}^{-\lambda t}\to0$ (croissances comparées) et
    $\cos$, $\sin$ sont bornées. Toute solution tend vers
    \res{0} : le régime libre s'éteint.
\end{questions}
\begin{center}
\begin{tikzpicture}[x=1.2cm,y=2.4cm,line cap=round]
  \draw[otRule,line width=0.3pt] (0,0.5) -- (6.2,0.5) (0,1) -- (6.2,1);
  \draw[otMuted,line width=0.5pt,-{Stealth[length=1.6mm]}] (0,0) -- (6.4,0);
  \draw[otMuted,line width=0.5pt,-{Stealth[length=1.6mm]}] (0,-0.2) -- (0,1.2);
  \draw[otNavy!55,line width=1pt] plot[domain=0:6.2,samples=120]
    ({\x},{1.5*exp(-\x)-0.5*exp(-3*\x)});
  \draw[otGreenDark,line width=1.1pt] plot[domain=0:6.2,samples=120]
    ({\x},{(1+\x)*exp(-\x)});
  \draw[otNavy,line width=1.2pt] plot[domain=0:6.2,samples=200]
    ({\x},{exp(-\x)*(cos(deg(sqrt(3)*\x))+sin(deg(sqrt(3)*\x))/sqrt(3))});
  \node[font=\scriptsize,text=otNavy!70,anchor=west] at (2.3,0.42) {apériodique};
  \node[font=\scriptsize,text=otGreenDark,anchor=west] at (1.6,0.68) {critique};
  \node[font=\scriptsize,text=otNavy,anchor=north west] at (1.9,-0.07) {pseudo-périodique};
  \node[font=\tiny,text=otMuted,anchor=east] at (0,1) {$1$};
  \node[font=\scriptsize,text=otMuted,below] at (6.3,0) {$t$};
\end{tikzpicture}\\[0.5mm]
\legende{Les trois régimes libres ($y(0)=1$, $y'(0)=0$) : le régime
critique est celui qui revient le plus vite à l'équilibre sans osciller.}
\end{center}
\end{corrige}
\end{blocLivre}

\begin{blocLivre}{cor-4-25b}
\begin{corrigeSuite}
\partieExo{B}{Le régime forcé ($f(t)=\cos\omega t$)}
\begin{questions}[start=4]
  \item Pour $z=Z\mathrm{e}^{\mathrm{i}\omega t}$ :
    $z''+2\lambda z'+\omega_0^{2}z
    =\bigl(-\omega^{2}+2\mathrm{i}\lambda\omega+\omega_0^{2}\bigr)Z\mathrm{e}^{\mathrm{i}\omega t}$,
    et le facteur est non nul (partie imaginaire $2\lambda\omega\neq0$) :
    $\res{Z=\dfrac{1}{\omega_0^{2}-\omega^{2}+2\mathrm{i}\lambda\omega}}$
    convient.
  \item Les coefficients étant réels, $y_p=\operatorname{Re}z$ est
    solution de $(E_f)$ avec $f=\operatorname{Re}\mathrm{e}^{\mathrm{i}\omega t}=\cos\omega t$.
    Écrivons $Z=A\,\mathrm{e}^{-\mathrm{i}\varphi}$ avec $A=|Z|=A(\omega)$ :
    $\mathrm{e}^{\mathrm{i}\varphi}=A\bigl(\omega_0^{2}-\omega^{2}+2\mathrm{i}\lambda\omega\bigr)$,
    soit
    \[
    \cos\varphi=(\omega_0^{2}-\omega^{2})\,A(\omega),\qquad
    \sin\varphi=2\lambda\omega\,A(\omega)>0 ,
    \]
    d'où $\varphi\in\intOO{0}{\pi}$, et
    $y_p(t)=\operatorname{Re}\bigl(A\mathrm{e}^{\mathrm{i}(\omega t-\varphi)}\bigr)
    =\res{A(\omega)\cos(\omega t-\varphi)}$.
  \item Par linéarité, $y-y_p$ est solution de $(E_0)$, donc tend vers
    $0$ (question 3) : $\res{y(t)-y_p(t)\to0}$. La réponse « oublie » les
    conditions initiales et ne garde que la pulsation imposée.
\end{questions}
\partieExo{C}{La résonance}
\begin{questions}[start=7]
  \item Avec $u=\omega^{2}$ :
    $D(u)=(\omega_0^{2}-u)^{2}+4\lambda^{2}u=u^{2}-2(\omega_0^{2}-2\lambda^{2})u+\omega_0^{4}$,
    trinôme minimal en $u^{*}=\omega_0^{2}-2\lambda^{2}$. Si
    $\lambda<\frac{\omega_0}{\sqrt2}$, $u^{*}>0$ : $D$ décroît puis croît
    sur $\R_{+}$, et $A=D^{-1/2}$ est maximale en
    $\res{\omega_r=\sqrt{\omega_0^{2}-2\lambda^{2}}}$, avec
    $D(u^{*})=\omega_0^{4}-(\omega_0^{2}-2\lambda^{2})^{2}
    =4\lambda^{2}(\omega_0^{2}-\lambda^{2})$, d'où
    $\res{A_{\max}=\dfrac{1}{2\lambda\sqrt{\omega_0^{2}-\lambda^{2}}}}$.
    Si $\lambda\geqslant\frac{\omega_0}{\sqrt2}$, $u^{*}\leqslant0$ : $D$
    croît sur $\R_{+}$, $A$ décroît, pas de résonance.
  \item $A(0)=\dfrac{1}{\omega_0^{2}}$ et $A(\omega_0)=\dfrac{1}{2\lambda\omega_0}$,
    d'où $\res{\dfrac{A(\omega_0)}{A(0)}=\dfrac{\omega_0}{2\lambda}}$.
    Quand $\lambda\to0$ : $\omega_r\to\omega_0$ et
    $A_{\max}\sim\dfrac{1}{2\lambda\omega_0}\to+\infty$. Un système peu
    amorti, excité près de sa pulsation propre, répond avec une amplitude
    considérable : c'est pourquoi on amortit les structures (ponts,
    machines, suspensions) soumises à des sollicitations périodiques. La
    figure de l'énoncé montre le pic s'affiner et monter quand $\lambda$
    diminue, et disparaître pour $\lambda=0{,}8>\frac{1}{\sqrt2}$.
\end{questions}
\end{corrigeSuite}
\end{blocLivre}
